\begin{pastexam}[label={px:ncu-113-dsa-20}]{中央資工所 113 學年度 第 20 題}
（複選；每題 5 分，全部答對才給分，答錯倒扣一分）

Which of the following statements describe the differences between Dynamic Programming and
Divide-and-Conquer?
\begin{enumerate}[label=(\Alph*)]
\item Whether the subproblems overlap or not
\item The division of problems
\item The combination of subproblems
\item The way the base case is solved
\item The depth of recurrence
\end{enumerate}
\end{pastexam}

\begin{pastsoln}
答案：(A)。

這題的選項寫得很簡短，判斷的標準是：這一項是不是兩種方法本質上不同的地方。

兩種方法的骨架其實一樣：把問題切成較小的子問題（division）、遞迴解子問題、碰到最小的情況直接解（base
case）、再把子問題的答案組合起來（combination）。 動態規劃的遞迴式和分治的遞迴式寫起來也一樣是「大問題
用小問題表示」。 真正的差別在於子問題會不會重複出現（資結書第~18.1.1~節）：
\begin{itemize}
\item 分治的子問題互不重疊（合併排序的左右兩半沒有共同的部分），每個子問題只會被解一次，直接遞迴就好。
\item 動態規劃的子問題大量重疊（費氏數的 $F(n-2)$ 同時被 $F(n)$、$F(n-1)$ 用到），所以要把答案存起來，
用備忘錄或由下而上填表，避免重算。
\end{itemize}

逐項檢查：
\begin{itemize}
\item (A) 對，就是上面說的差別。
\item (B) 錯。 兩者都把問題切成子問題；切法依問題而定（切一半、切掉一個元素、列舉切點……），不是區分兩者
的依據。
\item (C) 錯。 兩者都要把子問題的答案組合起來（合併排序的合併、DP 遞迴式裡的 $\min$／$\max$）。
\item (D) 錯。 兩者的基底情況都是直接回答最小的子問題。
\item (E) 錯。 遞迴深度由問題怎麼切決定，例如分治的二分搜尋深度 $\lg n$、動態規劃的費氏數深度 $n$，但也有
深度 $n$ 的分治和深度 $\lg n$ 的動態規劃，不是兩者的差別。
\end{itemize}
\end{pastsoln}
